\section{Evaluation protocol}
\label{sec:evaluation}

\subsection{Partitioning}

Users are randomly assigned to train ($70\%$), validation ($15\%$) and test
($15\%$), and every review by a given user lands in exactly one partition
(\texttt{grouped\_user\_split}, Section~\ref{sec:data}). The split is
reconstructed and checked for user overlap on every call
(\texttt{GroupSplit.\_\_post\_init\_\_}), and
\texttt{tests/test\_features\_leakage.py} asserts that attempting to build a
split from row-sliced (rather than user-grouped) partitions raises.

\subsection{Metrics}

\paragraph{Mean absolute ordinal error.} With $y_i \in \{0, 1, 2\}$ the true
level and $\hat{y}_i$ the predicted level,
\begin{equation}
  \text{MAE}_{\text{ord}} = \frac{1}{n} \sum_{i=1}^{n} \lvert y_i - \hat{y}_i \rvert.
\end{equation}
Unlike plain accuracy, this distinguishes a \emph{small}-predicted-as-\emph{large}
error (distance 2) from a \emph{small}-predicted-as-\emph{fit} error
(distance 1) -- the entire reason to treat the problem as ordinal.

\paragraph{Accuracy within one level.} The fraction of predictions at most
one ordinal step from the truth, $\frac{1}{n}\sum_i
\mathbb{1}[\lvert y_i - \hat{y}_i\rvert \le 1]$ -- a tolerant metric that
still distinguishes a model that is occasionally off by one level from one
that confuses \emph{small} and \emph{large} outright.

\paragraph{Macro F1.} Unweighted mean of per-class F1, reported because
accuracy alone is a weak signal under the dataset's class skew (the majority
baseline attains high accuracy and within-one-level scores by predicting
"fit" for everyone, but a macro F1 near the single-class floor).

\paragraph{Conformal coverage and average set size.} Empirical fraction of
test examples whose true label is included in the conformal set
(Equation~\ref{eq:prediction-set}), and the mean set cardinality. Coverage
should sit near the target $1-\alpha$; set size is the cost of that
guarantee -- a trivial predictor could "achieve" perfect coverage by
returning every class every time (as the majority baseline does in
Table~\ref{tab:conformal}), which is why set size must always be read
alongside coverage, not instead of it.

\paragraph{What we do not report, and why.} NDCG@K, Recall@K and MAP are
absent from this report by design, not omission; see
Section~\ref{sec:why-not-ranking} for the argument.

\subsection{Model selection and test discipline}

Every model's hyperparameters ($\ell_2$ strength for both logistic models)
were selected by inspecting validation-partition metrics; the test partition
is touched only once, by \texttt{scripts/export\_results\_table.py}, to
produce Table~\ref{tab:results}. The conformal calibration set is the
validation partition, disjoint from both the training data the model's
\emph{parameters} were fit on and the test data its \emph{coverage} is
measured on -- using the same partition for fitting the model and
calibrating the conformal wrapper would violate the exchangeability
assumption the coverage guarantee (Equation~\ref{eq:coverage-guarantee})
depends on.

\subsection{Training/serving parity}

\texttt{tests/test\_serving.py} asserts every precomputed payload's predicted
class is a member of its own conformal set (the fitted model's own top
choice can never be excluded by its own calibrated threshold), that
class probabilities sum to one, that request ids are unique, and that the
full payload is JSON-serializable -- the same check the reference
BPR/SVD portfolio projects use to catch NumPy-scalar serialization failures
before they reach a Lambda deployment.
